% !TeX root = ../../main.tex
\section{Background and Related Work}
\label{sec:mobiwac:related}


\subsection{From place embeddings to per-visit embeddings}
\label{sec:mobiwac:related-embeddings}
A place embedding transforms a POI into a vector so that similar or nearby
places sit close together in the vector space.
Deep Graph Infomax (DGI)~\cite{velickovic2019deep}, a general self-supervised
method for graphs, learns such vectors on a place network linked by
similarity, time, and distance, training them to tell the real network from a
copy with shuffled node features.
Hierarchical Graph Infomax (HGI) \cite{huang2023hgi} adds geographic structure,
organizing the network as place, then region, then city. Both produce one vector per place, so two visits to the same coffee shop look
identical to the model.
Contextual check-in representations instead give each visit its own vector: CTLE
\cite{lin2021ctle}, the closest example, learns a vector per visit by masking and
reconstructing parts of a user's check-in sequence. Our
representation, Check2HGI, also works at the check-in level but through a
different construction. CTLE is a sequence model, a Transformer that reads the
check-in sequence itself; Check2HGI remains a network model, keeping the
hierarchical place-to-region-to-city network and the same infomax objective,
now extended one level deeper, from individual places down to individual
check-ins (Figure~\ref{fig:mobiwac:dataflow}). The order of a user's visits still enters
the network, as links between consecutive check-ins
(Section~\ref{sec:mobiwac:method-rep}).

The novelty is this specific combination: per-visit context inside a
hierarchical graph-infomax representation. We compare against CTLE and a feature-concatenation control to separate the
combination from contextualization alone and from feature injection.

\subsection{Predicting the next category and the next region}
\label{sec:mobiwac:related-tasks}
A large body of work predicts the next place that a user will visit, the exact POI,
with recurrent and attention models such as DeepMove \cite{feng2018deepmove}, STAN
\cite{luo2021stan}, and GETNext \cite{yang2022getnext}. In contrast, we target two properties of the next visit, its category and its
region, rather than the exact place. We choose them because they are
what a mobility-aware service can act on, not to make the task easier; the region
alone spans hundreds to several thousand classes. Predicting
over a partition of the map is also the standard formulation in the human-mobility
literature, with a grid cell as the target~\cite{luca2021mobilitysurvey}; our
next-region task substitutes official neighborhood-scale units for grid cells. Our earlier
work~\cite{silva2025mtlnet} paired static category classification with next-category prediction
and found no consistent multitask gain; this chapter introduces the next-region task and adds the check-in-level
representation, on which sharing helps instead of hurting (Section~\ref{sec:mobiwac:results-part2}).

The field
increasingly models several granularities at once; in those systems, category and
region are auxiliary signals that help a primary next-place task
(MCMG \cite{sun2024mcmg}, HMT-GRN \cite{Lim2022}).
We study the pair as the object itself. To our knowledge, fine-grained region as an
end target of equal standing, rather than an auxiliary coarse grid cell, is
underexplored.
The nearest exceptions do not study our exact pairing:
DRRGNN~\cite{zhu2022drrgnn} forecasts a person's next activity region jointly
with a mobility-intention label, but over regions discovered per person rather
than a fixed citywide partition; a generative recommender~\cite{sun2025kgtb}
predicts category and region together, but only as auxiliary steps toward its
next-place ranking.

\subsection{Multitask learning for mobility}
\label{sec:mobiwac:related-mtl}
MCARNN \cite{Liao2018} jointly predicts activity and
location, and a recent hierarchy-aware model continues the
pattern \cite{wang2025hamtl}; in both, the location target is
the exact place. The closest alternative to our setting is
the cascade: predict the category first, then the place given the
category \cite{ye2013nextmove}. CSLSL
\cite{huang2024cslsl} is its current form, predicting in a chain (when, then what,
then where), with location as the primary output and category as an
intermediate step; CatDM \cite{yu2020catdm} likewise uses the predicted
category only to filter candidate places. We instead predict
category and region in parallel, as end targets of equal standing in one model
with a single forward pass, and we drop the next-place target entirely, so
neither task is an intermediate step toward a third. Since the cascade is the pattern that the field uses, we test the
choice directly (Section~\ref{sec:mobiwac:results-part2}). On optimization, we are conservative by design: joint
training with a fixed loss weighting is standard practice
\cite{caruana1997multitask}, not itself our contribution.
Reports find that gradient-balancing methods, which re-weight
the tasks' updates during training (PCGrad \cite{yu2020pcgrad}, Nash-MTL \cite{nash}), rarely improve on a
well-tuned fixed weighting with two tasks \cite{xin2022domtl}.
We confirm this at scale, screening nineteen loss and gradient balancers, including
the two methods named above, at their default configurations and a single seed on two
datasets, Alabama and Florida: none improved on a tuned fixed task weighting across both
tasks and both datasets. Two exceed equal weighting on next-category at Alabama,
Nash-MTL by $0.68$ points and scale normalization by $0.19$. At Florida, Nash-MTL falls
back below equal weighting on both tasks, while scale normalization still exceeds equal
weighting on next-category but collapses on next-region.
The reason is visible in the gradients. Measured during development on the same
joint architecture (on an earlier preparation of the data), the cosine
similarity between the next-category and next-region updates on the shared
trunk averages $+0.001$ across training (four seeds each on four Gowalla
states: Alabama, Arizona and Florida, which are three of the five U.S. datasets reported
here, and Georgia, which is not among the datasets this study reports; the largest per-dataset mean in
absolute value is $+0.0032$). A balancer therefore has no
conflict to resolve, consistent with the mechanism test in
Section~\ref{sec:mobiwac:results-part2}; this is a finding for this pair of tasks, not a
general rule. The same quantity, measured on the final model across seven datasets, is
positive at every one of them, so the absence of conflict is not an artifact of this
smaller set of runs.


\begin{figure}[htbp]
\centering
\includegraphics[width=\textwidth]{figures/mobiwac/fig1_dataflow.pdf}
\caption{From a check-in sequence to two predictions: each visit is embedded in a four-level graph
(check-in, place, region, city) trained without task labels. The graph yields two sets of
vectors, one per visit and one per region; sliding nine-visit windows of each feed a single
model that outputs the next category and the next region.}
\label{fig:mobiwac:dataflow}
\end{figure}
